\chapter{Single Fields}

\section{Electrostatics}
\subsection{Results}
The electric field intensity vector $\bE$ can be derived from the
nodal degree of freedom, the electric potential $V$ by applying a
gradient operator. This element quantity is calculated by applying
the config-file-commands

\begin{Verbatim}[commandchars=@\[\]]
@PYaZ[<storeResults]@PYaZ[>]
   @PYaZ[<elemResult] @PYaQ[type=]@PYaB["elecFieldIntensity"] @PYaQ[region=]@PYaB["elecst2d"]@PYaZ[/>]
@PYaZ[</storeResults>]
\end{Verbatim}


In the command  \xat{region}, all areas in which the electric field has to
be calculated, must be listed. If no region  is specified, the elecric
field will be calculated in all regions where the PDE is defined. To
calculate the electric energy, one has to use the command 

\begin{Verbatim}[commandchars=@\[\]]
      
@PYaZ[<storeResults]@PYaZ[>]
   @PYaZ[<elemResult] @PYaQ[type=]@PYaB["elecEnergy"] @PYaQ[region=]@PYaB["reg2"]@PYaZ[/>]
@PYaZ[</storeResults>]
\end{Verbatim}

in the config-file. The result is written to the info-file.


\subsection{Example}

{\footnotesize
\begin{Verbatim}[commandchars=@\[\]]
@PYaN[<?xml version="1.0"?>]

@PYaZ[<cfsSimulation] @PYaQ[xmlns=]@PYaB["http://www.cfs++.org"]@PYaZ[>]

   @PYaZ[<fileFormats]@PYaZ[>]
    @PYaZ[<output]@PYaZ[>]
      @PYaZ[<unv]@PYaZ[/>]
    @PYaZ[</output>]
    @PYaZ[<materialData] @PYaQ[file=]@PYaB["mat.xml"] @PYaQ[format=]@PYaB["xml"]@PYaZ[/>]
   @PYaZ[</fileFormats>]

  @PYaZ[<domain] @PYaQ[geometryType=]@PYaB["axi"]@PYaZ[>]
   @PYaZ[<regionList]@PYaZ[>]
    @PYaZ[<region] @PYaQ[name=]@PYaB["gndEl"]  @PYaQ[material=]@PYaB["steel"]@PYaZ[/>]
    @PYaZ[<region] @PYaQ[name=]@PYaB["hotEl"]  @PYaQ[material=]@PYaB["steel"]@PYaZ[/>]
    @PYaZ[<region] @PYaQ[name=]@PYaB["dielEl"] @PYaQ[material=]@PYaB["silicon"]@PYaZ[/>]
    @PYaZ[<region] @PYaQ[name=]@PYaB["airEl"]  @PYaQ[material=]@PYaB["air"]@PYaZ[/>]
   @PYaZ[</regionList>]

   @PYaZ[<nodeList]@PYaZ[>]
    @PYaZ[<nodes] @PYaQ[name=]@PYaB["hot"]@PYaZ[/>]
    @PYaZ[<nodes] @PYaQ[name=]@PYaB["gnd"]@PYaZ[/>]
   @PYaZ[</nodeList>]
  @PYaZ[</domain>]

 @PYaZ[<sequenceStep]@PYaZ[>]
  @PYaZ[<analysis]@PYaZ[>]
   @PYaZ[<static]@PYaZ[/>]
  @PYaZ[</analysis>]

  @PYaZ[<pdeList]@PYaZ[>]
    @PYaZ[<electrostatic]@PYaZ[>]
     @PYaZ[<regionList]@PYaZ[>]
       @PYaZ[<region] @PYaQ[name=]@PYaB["gndEl"]@PYaZ[/>]
       @PYaZ[<region] @PYaQ[name=]@PYaB["hotEl"]@PYaZ[/>]
       @PYaZ[<region] @PYaQ[name=]@PYaB["dielEl"]@PYaZ[/>]
       @PYaZ[<region] @PYaQ[name=]@PYaB["airEl"]@PYaZ[/>]
     @PYaZ[</regionList>]

      @PYaZ[<bcsAndLoads]@PYaZ[>]
        @PYaZ[<dirichletHom] @PYaQ[name=]@PYaB["gnd"]@PYaZ[/>]
        @PYaZ[<dirichletInhom] @PYaQ[name=]@PYaB["hot"] @PYaQ[value=]@PYaB["1.0"]@PYaZ[/>]
      @PYaZ[</bcsAndLoads>]

      @PYaZ[<storeResults]@PYaZ[>]
        @PYaZ[<elemResult] @PYaQ[type=]@PYaB["elecFieldIntensity"]@PYaZ[>]
         @PYaZ[<allRegions]@PYaZ[/>]
        @PYaZ[</elemResult>]
        @PYaZ[<elemResult] @PYaQ[type=]@PYaB["elecEnergy"]@PYaZ[>]
         @PYaZ[<allRegions]@PYaZ[/>]
        @PYaZ[</elemResult>]
      @PYaZ[</storeResults>]

    @PYaZ[</electrostatic>]
  @PYaZ[</pdeList>]
 @PYaZ[</sequenceStep>]
@PYaZ[</cfsSimulation>]
\end{Verbatim}

}


\newpage
\section{Mechanics}
\subsection{Degree of freedoms}
In the 3D case we have {\bf x, y, z} and in the 2D case
(axisymmetric as well as plane) {\bf x, y}.


\subsection{Subtypes}

A specification of the geometric type of the problem is done on the top-level
via the \xel{domain} attribute \xat{geometryType}, which allows the
values \xatv{3d}, \xatv{plane}
and \xatv{axi} for its \xat{geometryType} attribute. For a mechanic
PDE this definition
must be repeated in the \xel{mechanic} element via its \xat{subtype}
entry.
\begin{Verbatim}[commandchars=@\[\]]
       @PYaZ[<mechanic] @PYaQ[subType=]@PYaB["stype"]@PYaZ[>]
\end{Verbatim}

\noindent Here stype can currently take one of the following three values

\begin{itemize}
\item planeStrain \emph{(the default)}
\item 3d
\item axi
\end{itemize}

\emph{Note:} The \xat{subType} value must of course match to the
\xat{geometryType}
value of the geometry specification and a mixture of different subTypes within
one analysis is thus not possible.

\subsection{Node to Ground Springs}
In mechanic, it is possible to consider node to ground springs. In a
static analysis one can consider a stiffness value and
in a transient analysis additionally a mass and a damping value.
With the parameter name an existing selection of nodes must be given
as it is the case for loads.
\begin{Verbatim}[commandchars=@\[\]]
@PYaZ[<bcsAndLoads]@PYaZ[>]
 @PYaZ[<spring] @PYaQ[name=]@PYaB["springNodes"] @PYaQ[dof=]@PYaB["y"] @PYaQ[massValue=]@PYaB["1"] @PYaQ[dampingValue=]@PYaB["1"] 
  @PYaQ[stiffnessValue=]@PYaB["1"]@PYaZ[/>]
@PYaZ[</bcsAndLoads>]
\end{Verbatim}


\subsection{Nodal Results}
Possible nodal results are

\begin{Verbatim}[commandchars=@\[\]]
 @PYaZ[<nodeResult] @PYaQ[type=]@PYaB["mechDisplacement"]@PYaZ[>]
  @PYaZ[<regionList]@PYaZ[>]
   @PYaZ[<region] @PYaQ[name=]@PYaB["bar"]@PYaZ[/>]
  @PYaZ[</regionList>]
 @PYaZ[</nodeResult>]
 @PYaZ[<nodeResult] @PYaQ[type=]@PYaB["mechVelocity"]@PYaZ[>]
  @PYaZ[<regionList]@PYaZ[>]
   @PYaZ[<region] @PYaQ[name=]@PYaB["bar"]@PYaZ[/>]
  @PYaZ[</regionList>]
 @PYaZ[</nodeResult>]
 @PYaZ[<nodeResult] @PYaQ[type=]@PYaB["mechAcceleration"]@PYaZ[>]
  @PYaZ[<regionList]@PYaZ[>]
   @PYaZ[<region] @PYaQ[name=]@PYaB["bar"]@PYaZ[/>]
  @PYaZ[</regionLit>]
 @PYaZ[</nodeResult>]
\end{Verbatim}


\subsection{Element Results}
The command for calculating the mechanical stresses is
\begin{Verbatim}[commandchars=@\[\]]
      
@PYaZ[<storeResults]@PYaZ[>]
   @PYaZ[<elemResult] @PYaQ[type=]@PYaB["mechStress"]@PYaZ[>]
    @PYaZ[<regionList]@PYaZ[>]
     @PYaZ[<region] @PYaQ[name=]@PYaB["bar"]@PYaZ[/>]
    @PYaZ[</regionList>]
   @PYaZ[</elemResult>]
@PYaZ[</storeResults>]
\end{Verbatim}


\subsection{Region Results}
There is one surface region result. The calculated value is of type history and
written to the \textit{info.xml} file or the filesystem if \xel{text} output
is enabled.

\subsubsection{Deformed Volume}
Given by 
\begin{Verbatim}[commandchars=@\[\]]
      
@PYaZ[<storeResults]@PYaZ[>]
  @PYaZ[<surfRegionResult] @PYaQ[type=]@PYaB["volumeAboveDefSurf"] @PYaQ[dof=]@PYaB["z"]@PYaZ[>]
    @PYaZ[<surfRegionList]@PYaZ[>]
      @PYaZ[<surfRegion] @PYaQ[name=]@PYaB[".."]@PYaZ[/>]
    @PYaZ[</surfRegionList>]
 @PYaZ[</surfRegionResult>]
@PYaZ[</storeResults>]
\end{Verbatim}

the displacements of the given direction are integrated over the given surface.
The result is for harmonic analysis a complex value. Of interest is them the
norm/length of the calculated vector.

\subsection{Example}

{\footnotesize 
\begin{Verbatim}[commandchars=@\[\]]
@PYaN[<?xml version="1.0"?>]
@PYaZ[<cfsSimulation] @PYaQ[xmlns=]@PYaB["http://www.cfs++.org"]@PYaZ[>]

 @PYaZ[<fileFormats]@PYaZ[>]
   @PYaZ[<output]@PYaZ[>]
    @PYaZ[<unv]@PYaZ[/>]
   @PYaZ[</output>]
   @PYaZ[<materialData] @PYaQ[file=]@PYaB["mat.xml"] @PYaQ[format=]@PYaB["xml"]@PYaZ[/>]
 @PYaZ[</fileFormats>]

 @PYaZ[<domain] @PYaQ[geometryType=]@PYaB["plane"]@PYaZ[>]
   @PYaZ[<regionList]@PYaZ[>]
    @PYaZ[<region] @PYaQ[name=]@PYaB["cantilev"]  @PYaQ[material=]@PYaB["steel"]@PYaZ[/>]
   @PYaZ[</regionList>]
   @PYaZ[<nodeList]@PYaZ[>]
    @PYaZ[<nodes] @PYaQ[name=]@PYaB["leftN"]@PYaZ[/>]
    @PYaZ[<nodes] @PYaQ[name=]@PYaB["rightN"]@PYaZ[/>]
   @PYaZ[</nodeList>]
 @PYaZ[</domain>]

 @PYaZ[<sequenceStep]@PYaZ[>]
   @PYaZ[<analysis]@PYaZ[>]
    @PYaZ[<static]@PYaZ[/>]
   @PYaZ[</analysis>]

   @PYaZ[<pdeList]@PYaZ[>]
    @PYaZ[<mechanic] @PYaQ[subtype=]@PYaB["planeStrain"]@PYaZ[>]
     @PYaZ[<regionList]@PYaZ[>]
      @PYaZ[<region] @PYaQ[name=]@PYaB["cantilev"]@PYaZ[/>]
     @PYaZ[</regionList>]
     @PYaZ[<bcsAndLoads]@PYaZ[>]
      @PYaZ[<dirichletHom] @PYaQ[name=]@PYaB["leftN"] @PYaQ[dof=]@PYaB["x"]@PYaZ[/>]
      @PYaZ[<dirichletHom] @PYaQ[name=]@PYaB["leftN"] @PYaQ[dof=]@PYaB["y"]@PYaZ[/>]
      @PYaZ[<dirichletHom] @PYaQ[name=]@PYaB["rightN"] @PYaQ[dof=]@PYaB["x"]@PYaZ[/>]
      @PYaZ[<load]         @PYaQ[name=]@PYaB["rightN"] @PYaQ[dof=]@PYaB["y"] @PYaQ[value=]@PYaB["-1.0"]@PYaZ[/>]
     @PYaZ[</bcsAndLoads>]
     @PYaZ[<storeResults]@PYaZ[>]
      @PYaZ[<nodeResult] @PYaQ[type=]@PYaB["mechDisplacement"]@PYaZ[>]
       @PYaZ[<allRegions]@PYaZ[/>]
      @PYaZ[</nodeResult>]
     @PYaZ[</storeResults>]
    @PYaZ[</mechanic>]
   @PYaZ[</pdeList>]
 @PYaZ[</sequenceStep>]

@PYaZ[</cfsSimulation>]
\end{Verbatim}

}


\newpage
\section{Acoustics}

\subsection{General Setup}
The \xel{acoustic} element is appended by the following child elements
\begin{itemize}
\item \xel{region}
\item \xel{surface}
\item \xel{timeSteppingFormulation} effective stiffness formulation 
is default(optional), effective mass matrix formulation can be chosen unless one
is not using the fractional damping model
\item \xel{bcsAndLoads} (optional)
\item \xel{nonLinear}\footnote{Other Nonlinearities than nonlinear acoustic
 PDE, which can be chosen by an attribute of \xel{region} element. This is
 currently not implemented.}
\item \xel{storeResults} (optional)
\end{itemize}
and the following attributes
\begin{itemize}
\item \xel{pressFormu} used to choose for an acoustic pressure PDE
formulatation \\
with \xatv{acouPressure} as value.
\item \xat{subType} for flownoise computations set this to
\xatv{combustionNoise} or \xatv{standard} 
\end{itemize}

\subsection{Formulation}
For linear acoustics, the PDE in pressure and acoustic potential has the same 
form. Allthough default is acoustic potential. If acoustics is coupled to 
other PDEs, acoustic potential is obligatory. Westervelt's nonlinear PDE is 
in pressure, Kuznetsov's nonlinear PDE is in acoustic potential.

\subsection{Damping and nonlinear acoustic simulation}
Damping and nonlinear acoustic PDE formulations can be chosen for each 
region individually. Damping is a child element of \xel{region},
type is either no damping, rayleigh damping (which is not physical), 
thermoviscous damping and fractional damping according to a 
frequency power law.
\xat{nonLinear} comes as attribute to region. It can take the value
\xatv{no},
\xatv{westervelt}, or \xatv{kuznetsov}.

\subsection{Boundary Conditions}
Unique to the acoustic PDE is the possibiliy to prescribe inhomogeneous Neumann
boundary conditions. In acoustic potential formulation this means a normal
velocity on a surface. Remember, that the specified surface elements have to
appear in the \xel{domain} section and in the \xel{acoustic} section
below the region definitions. The syntax is:

\begin{Verbatim}[commandchars=@\[\]]
 
@PYaZ[<bcsAndLoads]@PYaZ[>]
  @PYaZ[<neumannInhom] @PYaQ[name=]@PYaB["transducer"] @PYaQ[value=]@PYaB["1e-6"] @PYaQ[fileName=]@PYaB["sweep5M.dat"]@PYaZ[/>]
@PYaZ[</bcsAndLoads>]
\end{Verbatim}


\subsection{Nodal Results}
The config-file command for storing e.g. the acoustic potential in the result
file is

\begin{Verbatim}[commandchars=@\[\]]
      
@PYaZ[<storeResults]@PYaZ[>]
   @PYaZ[<nodeResult] @PYaQ[type=]@PYaB["acouPotential"]@PYaZ[>]
      @PYaZ[<regionList]@PYaZ[>]
        @PYaZ[<region] @PYaQ[name=]@PYaB["fluid1"]@PYaZ[/>]
      @PYaZ[</regionList>]
   @PYaZ[</nodeResult>] 
@PYaZ[</storeResults>]
\end{Verbatim}


Analogously one can specify storing of certain node values with the help of the
\xel{nodeList}, by the sub element \xel{nodes} and \xat{name}
attribute. 

\begin{Verbatim}[commandchars=@\[\]]
    
@PYaZ[<storeResults]@PYaZ[>]
  @PYaZ[<nodeResult] @PYaQ[type=]@PYaB["acouPotential"]@PYaZ[>]
    @PYaZ[<nodeList]@PYaZ[>]
     @PYaZ[<nodes] @PYaQ[name=]@PYaB["historyPoints"]@PYaZ[/>]
    @PYaZ[</nodeList>]
  @PYaZ[</nodeResult>] 
@PYaZ[</storeResults>]
\end{Verbatim}


Remember, that the quantity, you want to store as result, depends on the
formulation of the PDE. It is straightforward to retrieve the quantity acoustic
potential, when the formulation of the PDE is in acoustic potential. But if your
formulation is in acoustic pressure, it is nonsense to retrieve the quantity
acoustic potential, because this would need integration. A combination, which
does work, is formulation in acoustic potential and results in acoustic
pressure. This will be discussed in the next section.

\subsection{Element Result}

Because $p = \rho \partial \psi / \partial t$, and we can only assign a certain
density to an element in the computational domain, pressure is an element result
for acoustic potential formulation. It is specified with

\begin{Verbatim}[commandchars=@\[\]]
@PYaZ[<storeResults]@PYaZ[>]
  @PYaZ[<elemResult] @PYaQ[type=]@PYaB["acouPressure"]@PYaZ[>]
   @PYaZ[<elemList]@PYaZ[>]
    @PYaZ[<elems] @PYaQ[name=]@PYaB["receiver"]@PYaZ[/>]
   @PYaZ[</elemList>]
  @PYaZ[</elemResult>]
@PYaZ[</storeResults>]
\end{Verbatim}



\subsection{Region Result}
There are two surface region results. The result is of type history and  written
to the \textit{info.xml} file or to the filesystem if \xat{text} output is
enabled.   

\begin{Verbatim}[commandchars=@\[\]]
@PYaZ[<storeResults]@PYaZ[>]
  @PYaZ[<surfRegionResult] @PYaQ[type=]@PYaB["acouPower|acouForce"]@PYaZ[>]
    @PYaZ[<surfRegionList]@PYaZ[>]
      @PYaZ[<surfRegion] @PYaQ[name=]@PYaB[".."] @PYaZ[/>]
    @PYaZ[</surfRegionList>]
  @PYaZ[</surfRegionResult>]
@PYaZ[</storeResults>]
\end{Verbatim}


\subsubsection{Acoustic Power}
The acoustic power calculates
\begin{equation}
  P = \frac{1}{2} \int v\, p^*\,d\Gamma \nonumber
\end{equation}
with $v$ the complex acoustic particle velocity and $p^*$ the complex conjugate
acoustic pressure. Of special interest is the acoustic effective power
\begin{equation}
  P_{ac} = \frac{1}{2} \int {\rm Re} \lbrace p\, v^* \rbrace \,d\Gamma \nonumber
\end{equation}
which coincides with the real part of the calculated $P$.


\subsection{Absorbing Boundary Conditions (ABC)} 
In the config file, the command

\begin{Verbatim}[commandchars=@\[\]]
    
@PYaZ[<bcsAndLoads]@PYaZ[>]  
   @PYaZ[<absorbingBCs] @PYaQ[name=]@PYaB["abc"]@PYaZ[/>]
@PYaZ[</bcsAndLoads>]
\end{Verbatim}

      
is needed, which defines the boundaries at which the \xatv{abc} are
working.
It should be clear, that the dimension of the boundary elements are
one order lower compared to the dimension of the "domain" elements.
Therefore, the boundary elements have to be listed with the config-file
command

\begin{Verbatim}[commandchars=@\[\]]
@PYaZ[<domain]@PYaZ[>]
  @PYaZ[<surfRegionList]@PYaZ[>]
   @PYaZ[<surfRegion] @PYaQ[name=]@PYaB["abc"]@PYaZ[/>]
  @PYaZ[</surfRegionList>]
@PYaZ[</domain>]
\end{Verbatim}


\subsection{Aeroacoustics}
In aeroacoustics, sound sources are computed from a CFD simulation. This is
done by another tool, called \texttt{cplreader} (see the \texttt{cplreader}
manual for details). The sound sources are then used as a load on the right
hand side of the acoustic PDE. Therefore the corresponding XML element is
called \xel{rhsValues}. Here is an example:

\begin{Verbatim}[commandchars=@\[\]]
@PYaZ[<bcsAndLoads]@PYaZ[>]
 @PYaZ[<rhsValues] @PYaQ[region=]@PYaB["fluid"] @PYaQ[inputId=]@PYaB["acouGrid"]@PYaZ[>]
  @PYaZ[<factor]@PYaZ[>]1.0@PYaZ[</factor>]
  @PYaZ[<asyncSteps]@PYaZ[>]nearest@PYaZ[</asyncSteps>]
  @PYaZ[<interpolation] @PYaQ[justInterpolate=]@PYaB["yes"]@PYaZ[>]
    @PYaZ[<srcRegions] @PYaQ[names=]@PYaB["partition1"] @PYaQ[inputId=]@PYaB["sourceGrid"] @PYaQ[coordSysId=]@PYaB["default"]@PYaZ[/>]
    @PYaZ[<xyPlane] @PYaQ[z=]@PYaB["0.0"] @PYaQ[tol=]@PYaB["1.0e-10"]@PYaZ[/>]
    @PYaZ[<tolerances]@PYaZ[>]
      @PYaZ[<global] @PYaQ[start=]@PYaB["1e-4"] @PYaQ[end=]@PYaB["1e-3"] @PYaQ[inc=]@PYaB["9e-4"]@PYaZ[/>]
      @PYaZ[<local] @PYaQ[start=]@PYaB["1e-3"] @PYaQ[end=]@PYaB["5.1e-2"] @PYaQ[inc=]@PYaB["1e-2"]@PYaZ[/>]
    @PYaZ[</tolerances>]
    @PYaZ[<restartFileMode]@PYaZ[>]rw@PYaZ[</restartFileMode>]
    @PYaZ[<nodeWarnings] @PYaQ[display=]@PYaB["verbose"] @PYaQ[writeNodes=]@PYaB["yes"]@PYaZ[/>]
  @PYaZ[</interpolation>]
 @PYaZ[</rhsValues>]
@PYaZ[</bcsAndLoads>]
\end{Verbatim}


The attribute \xat{region} identifies the region on which the source term
is to be applied; it is always required. If there is no sub-element
\xel{interpolation} then source terms are read from the file given
through
the attribute \xat{inputId} (which defaults to \xatv{default}).

Inside the \xel{factor} element one can supply a mathematical expression
(see chapter~\ref{chap:MathExpr}) which is multiplied with the source terms.
The expression may depend on time or frequency, and space. This is especially
useful in a transient simulation that begins after the flow field has had some
time to develop turbulence. In this case, immediately switching on the full
source term can lead to numerical errors. A ramp function has proven to give
much better results. A commonly used ramp is
\verb|if( t lt T; sin(2*pi*t/(4*T))^2; 1)|, where \texttt{T} must be replaced
with the time duration of the ramp.

The \xel{asyncSteps} element tackles that fact that time steps in CFD are
usually smaller than in computational acoustics. Setting this value to
\xelv{nearest} selects source terms from the time step that is nearest to
the
one that is currently being computed. \xel{interpolate} will result in a
linear temporal interpolation of source terms between the two time steps that
the current time lies in between. In order to use the default behavior of
CFS++, selecting data by step number, set attribute \xat{justInterpolate}
to attribute value \xatv{no} here.

The \xel{interpolation} element tells CFS++ to interpolate source terms from
a different grid. This is required, if a CFD grid is not suitable for
acoustics. The attribute \xat{justInterpolate} may be set to
\xatv{yes} if
just the interpolation should be computed, but not the acoustic field. The
\xel{srcRegions} element defines the CFD grid; \xat{names} is the list
of regions (separated by space, \verb|;| or \verb#|#) that contain
source terms; \xat{inputId} identifies the input file from which to read
the
CFD grid and source terms; optionally \xat{coordSysId} may specify the
coordinate system of the CFD grid if it is different from the acoustics grid.

The \xel{xyPlane} element is only relevant if the acoustics grid is
two-dimensional, but the CFD grid is three-dimensional. Then one may specify
the \xat{z} coordinate of the x-y-plane (defaults to 0). Additionally
\xat{tol} gives the tolerance of \xat{z} in meters, so that nodes in
the
CFD grid need not lie exactly on the plane.

Inside the \xel{tolerances} element one may specify a range for the global
and local tolerances, which are used to determine the interpolation weights.
The global tolerance (in meters) determines the space around a node in the CFD
grid, when its surrounding element in the acoustics grid is searched for. The
local tolerance is the space around an element (in the element's local
coordinate system) where nodes are still accepted, although they lie outside of
the element. These tolerance are most relevant for domains with curved
boundaries, which are discretized very differently by CFD grid and acoustics
grid. Tolerance ranges are swept in order to use small tolerances wherever
possible and large tolerances only in places where it is necessary.

\xel{restartFileMode} determines if a restart file containing the
interpolation weights is to be read (\xelv{r}), written (\xelv{w}) or
both
(\xelv{rw}). In this way the interpolation weights need to be computed only
once and can be reused if a simulation is restarted.

The \xel{nodeWarnings} element defines how nodes are handled which could not
be mapped during interpolation. Any message is suppressed if the attribute
\xat{display} is set to \xatv{no}. The total number of missed nodes
is
reported, when set to \xatv{yes}. A message for each single node is
generated
by a value of \xatv{verbose}. Optionally a text file containing the node
numbers is written, if the attribute \xat{writeNodes} is set to
\xatv{yes}.

\subsection{Example}

{\footnotesize 
\begin{Verbatim}[commandchars=@\[\]]
@PYaN[<?xml version="1.0"?>]
@PYaZ[<cfsSimulation] @PYaQ[xmlns=]@PYaB["http://www.cfs++.org"]@PYaZ[>]

  @PYaZ[<fileFormats]@PYaZ[>]
    @PYaZ[<output]@PYaZ[>]
     @PYaZ[<gid]@PYaZ[/>]
    @PYaZ[</output>]
    @PYaZ[<materialData] @PYaQ[file=]@PYaB["mat.xml"] @PYaQ[format=]@PYaB["xml"]@PYaZ[/>]
  @PYaZ[</fileFormats>]

  @PYaZ[<domain] @PYaQ[geometryType=]@PYaB["plane"]@PYaZ[>]
    @PYaZ[<regionList]@PYaZ[>]
      @PYaZ[<region] @PYaQ[material=]@PYaB["softTissue"] @PYaQ[name=]@PYaB["reg1"]@PYaZ[/>]
      @PYaZ[<region] @PYaQ[material=]@PYaB["bone"] @PYaQ[name=]@PYaB["reg2"]@PYaZ[/>]
    @PYaZ[</regionList>]
    @PYaZ[<nodeList]@PYaZ[>]
      @PYaZ[<nodes] @PYaQ[name=]@PYaB["vp-restr"]@PYaZ[/>]
    @PYaZ[</nodeList>]
  @PYaZ[</domain>]

 @PYaZ[<sequenceStep]@PYaZ[>]

  @PYaZ[<analysis]@PYaZ[>]
    @PYaZ[<transient]@PYaZ[>]
      @PYaZ[<numSteps]@PYaZ[>]    600     @PYaZ[</numSteps>]
      @PYaZ[<firstDt]@PYaZ[>]     2.5e-08 @PYaZ[</firstDt>]
      @PYaZ[<stepSaveBeg]@PYaZ[>]   1     @PYaZ[</stepSaveBeg>]
      @PYaZ[<stepSaveEnd]@PYaZ[>] 600     @PYaZ[</stepSaveEnd>]
      @PYaZ[<stepSaveInc]@PYaZ[>]   1     @PYaZ[</stepSaveInc>]
      @PYaZ[<timeDataFile] @PYaQ[name=]@PYaB["sin1M40cos.dat"]@PYaZ[/>]
    @PYaZ[</transient>] 
  @PYaZ[</analysis>]

  @PYaZ[<pdeList]@PYaZ[>]
    @PYaZ[<acoustic] @PYaQ[formulation=]@PYaB["acouPotential"]@PYaZ[>]
      @PYaZ[<regionList]@PYaZ[>]
       @PYaZ[<region] @PYaQ[name=]@PYaB["reg1"]@PYaZ[/>]
       @PYaZ[<region] @PYaQ[name=]@PYaB["reg2"]@PYaZ[/>]
      @PYaZ[</regionList>]
      @PYaZ[<bcsAndLoads]@PYaZ[>]
        @PYaZ[<dirichletInhom] @PYaQ[fileName=]@PYaB["sin1M40cos.dat"] @PYaQ[name=]@PYaB["vp-restr"] @PYaQ[value=]@PYaB["1.0"]@PYaZ[/>]
      @PYaZ[</bcsAndLoads>]
      @PYaZ[<storeResults]@PYaZ[>]
        @PYaZ[<nodeResult] @PYaQ[type=]@PYaB["acouPotential"]@PYaZ[>]
          @PYaZ[<nodeList]@PYaZ[>]
            @PYaZ[<nodes] @PYaQ[name=]@PYaB["historyPoints"]@PYaZ[/>]
          @PYaZ[</nodeList>]
        @PYaZ[</nodeResult>]
      @PYaZ[</storeResults>]
    @PYaZ[</acoustic>]
  @PYaZ[</pdeList>]

 @PYaZ[</sequenceStep>]
@PYaZ[</cfsSimulation>]
\end{Verbatim}

}


\newpage
\section{Magnetics}
\subsubsection{Element Results}
In the case of magnetic field calculation, many element results can be
deduced by a postprocessing step:

\begin{itemize}
\item Magnetic induction (magnetic flux density) $\bB$
    \newline Derived from the magnetic vector potential $\bA$ by
    applying the curl operator.
  
\begin{Verbatim}[commandchars=@\[\]]
@PYaZ[<storeResults]@PYaZ[>]
   @PYaZ[<nodeResult] @PYaQ[type=]@PYaB["magPotential"]@PYaZ[>]
     @PYaZ[<regionList]@PYaZ[>]
      @PYaZ[<region] @PYaQ[name=]@PYaB["yoke"]@PYaZ[/>]
     @PYaZ[</regionList>]
   @PYaZ[</nodeResult>]
   @PYaZ[<elemResult] @PYaQ[type=]@PYaB["magFluxDensity"] @PYaZ[>]
     @PYaZ[<regionList]@PYaZ[>]
      @PYaZ[<region] @PYaQ[name=]@PYaB["yoke"]@PYaZ[/>]
     @PYaZ[</regionList>]
   @PYaZ[</elemResult>]
@PYaZ[</storeResults>]
\end{Verbatim}

  
\item Eddy current density $\bJ_{\rm e}$ 
    \newline Currents induced by a magnetic field in a conductive
    material. 

\begin{Verbatim}[commandchars=@\[\]]
@PYaZ[<storeResults]@PYaZ[>]
   @PYaZ[<elemResult] @PYaQ[type=]@PYaB["magEddyCurrent"]@PYaZ[>]
     @PYaZ[<regionList]@PYaZ[>]
      @PYaZ[<region] @PYaQ[name=]@PYaB["yoke"]@PYaZ[/>]
     @PYaZ[</regionList>]
   @PYaZ[</elemResult>]
@PYaZ[</storeResults>]
\end{Verbatim}

  
  
\item Magnetic energy $W$ 
    \newline Measure for the energy stored in the magnetic assembly.
  
\begin{Verbatim}[commandchars=@\[\]]
@PYaZ[<storeResults]@PYaZ[>]
   @PYaZ[<regionResult] @PYaQ[type=]@PYaB["magEnergy"]@PYaZ[>]
     @PYaZ[<regionList]@PYaZ[>]
      @PYaZ[<region] @PYaQ[name=]@PYaB["yoke"]@PYaZ[/>]
     @PYaZ[</regionList>]
   @PYaZ[</regionResult>]
@PYaZ[</storeResults>]
\end{Verbatim}


\item  Inductivity $L$ 

\begin{Verbatim}[commandchars=@\[\]]
@PYaZ[<coils]@PYaZ[>]
  @PYaZ[<currentCoil2d] @PYaQ[name=]@PYaB["..."]@PYaZ[>]
     ...
    @PYaZ[<saveFileL]@PYaZ[>] lFile.dat @PYaZ[</saveFileL>]
  @PYaZ[</currentCoil2d>]
@PYaZ[</coils>]
\end{Verbatim}


The inductivity $L$ of the magnetic assembly, as well as
the magnetic flux $\Phi$ are calculated and stored in a file with
the name "lFile.dat"  given in column oriented form:

{\em time inductivity magFlux}.  
  


\end{itemize}

\subsection{Coil Definition}

The list of coils has to be defined inside the tags

\begin{Verbatim}[commandchars=@\[\]]
@PYaZ[<coils]@PYaZ[>]

@PYaZ[</coils>]
\end{Verbatim}


\subsubsection{Coil Types}

In our simulation, we distinguish different coil types:
\begin{itemize}
\item Measurement coils
  \begin{itemize}
  \item 2D: \verb! <measurementCoil2d name="coil1">!
  \item 3D: \verb! <measurementCoil3d name="coil2">!
  \end{itemize}
\item Volage driven coils 
  \begin{itemize}
  \item 2D: \verb! <voltageCoil2d name="coil3">!
  \item 3D: \verb! <voltageCoil3d name="coil4">!
  \end{itemize}
\item Current driven coils 
  \begin{itemize}
  \item 2D: \verb! <currentCoil2d name="coil5">!
  \item 3D: \verb! <currentCoil3d name="coil6">!
  \end{itemize}
\end{itemize}

There are several parameters which must be defined for a coil. Every
coil needs an attribute {\xat{name}="\xatv{xyz}"}, standing nearby
the
definition of the coil:

\verb! <voltageCoil2d name="coil3">!

\noindent
Further on, the cross section of a winding has to be
specified for every coil:

\verb!<windingCrossSection> 1e-6 </windingCrossSection>!

\noindent
The number of windings $N$ results from the cross section of the wire. It
is calculated by the overall cross section of the coil $A_c$ devided
by the cross section of a winding $A_w$

\begin{equation}
  \label{eq:N}
  N=\frac{A_c}{A_w}
\end{equation}

\noindent
Depending on the coil type, several additional tags are needed. They
are described in the next sections.


\subsubsection{Measurement Coils}

\begin{itemize}
\item Coil voltage $U$ respectively coil current $I$

\begin{Verbatim}[commandchars=@\[\]]
@PYaZ[<measurementCoil2d] @PYaQ[name=]@PYaB["coil1"]@PYaZ[>]
  @PYaZ[<windingCrossSection]@PYaZ[>] 1e-6 @PYaZ[</windingCrossSection>]
   @PYaZ[<saveFileU]@PYaZ[>]  uiFile.dat  @PYaZ[</saveFileU>] 
  @PYaZ[<id]@PYaZ[>]1@PYaZ[</id>]
@PYaZ[</measurementCoil2d>]
\end{Verbatim}


   With this tag the voltage respectively the current of a
   special coil element can be calculated. It's stored in a file
   named {\em resultFile} in the column oriented form:

   {\em time voltage}.

\end{itemize}

\subsubsection{Current- and Voltage Driven Coils}
In this coil type, the amplitude, phase (in case of harmonic analysis)
and the input function (in case of dynamic analyses) of the driving
current respectively voltage has to be given:

\begin{itemize}
\item Transient case
\begin{Verbatim}[commandchars=@\[\]]
@PYaZ[<value]@PYaZ[>]       10           @PYaZ[</value>]
@PYaZ[<dynamics]@PYaZ[>]    timefile.dat @PYaZ[</dynamics>]
\end{Verbatim}


\item Harmonic case
\begin{Verbatim}[commandchars=@\[\]]
@PYaZ[<value]@PYaZ[>]       10           @PYaZ[</value>]
@PYaZ[<phase]@PYaZ[>]       90           @PYaZ[</phase>]
\end{Verbatim}

\end{itemize}

\noindent
The name of the time function as well as the sepcification of the
phase is optional. If it's not given, a constant current with an
amplitude of {\xel{value}} is assumed.  
Voltage driven coils additionally require the total ohmic resistance
of the coil for solving the network equation:

\verb!<resistance>  10 </resistance>!

\subsubsection{2D Coils}
For all types of 2D-coils, an identification number is
necessary:

\verb! <id>  2  <id>!

\noindent
Due to the 2D- or axisymmetric modeling, only the cross-section of the
coil is modeled. The coil-id is
used to identify coil cross-sections belonging together (fed by the
same current). 
The sign of the id also specifies the direction of the current.

\subsubsection{3D Current Driven Coils}
For 3D current driven coils, the direction of the current has to be
given explicitly. In case of a straight current flow, three
possibilities can be given:

\begin{Verbatim}[commandchars=@\[\]]
@PYaZ[<flowDirection]@PYaZ[>]
  @PYaZ[<direction] @PYaQ[dof=]@PYaB["x"] @PYaQ[value=]@PYaB["1"]@PYaZ[/>]
@PYaZ[</flowDirection>]!
\end{Verbatim}


\noindent
with {\xel direction} standing for one of 
\begin{itemize}
\item x
\item y
\item z
\item phi
\item r
\end{itemize}
   


\subsubsection{Coil Definition Examples}

The full description of a voltage driven coil in 2D with a winding
cross section of $10^{-3}\, \rm m^2$, a voltage excitation of 10\,V
and 10\,$\Omega$ resistance could be given as
\begin{Verbatim}[commandchars=@\[\]]
@PYaZ[<voltageCoil2d] @PYaQ[name=]@PYaB["coil1"]@PYaZ[>]
   @PYaZ[<windingCrossSection]@PYaZ[>] 0.001  @PYaZ[</windingCrossSection>]
   @PYaZ[<value]@PYaZ[>]                  10  @PYaZ[</value>]
   @PYaZ[<resistance]@PYaZ[>]             10  @PYaZ[</resistance>]
   @PYaZ[<id]@PYaZ[>]                      2  @PYaZ[</id>]
@PYaZ[</voltageCoil2d>]
\end{Verbatim}



\subsection{Permanent Magnets}
The use of permanent magnets in the simulation is very simple. In the
magnetic section, the command

\begin{Verbatim}[commandchars=@\[\]]
@PYaZ[<magnetic]@PYaZ[>]
...
  @PYaZ[<magnets]@PYaZ[>]
   @PYaZ[<magnet] @PYaQ[name=]@PYaB["perm"] @PYaQ[orientX=]@PYaB["0"] @PYaQ[orientY=]@PYaB["0"] @PYaQ[orientZ=]@PYaB["1"] @PYaZ[/>]
  @PYaZ[</magnets>]
...
@PYaZ[</magnetic>]
\end{Verbatim}

has to list all domains, which hold a permanent magnet. The
orientation and strength of the permanent magnetic field is given by
the magnetization vector $\bM$ given in Tesla described by the three
values {\xat{orientX} }, {\xat{orientY} } and {\xat{orientZ} }

\begin{equation}
  \label{eq:M}
  \bM = \left(
    \begin{tabular}{r}
      orientX\\
      orientY\\
      orientZ
    \end{tabular}
\right) \; .
\end{equation}

\subsubsection{Example}
{\footnotesize
\begin{Verbatim}[commandchars=@\[\]]
@PYaN[<?xml version="1.0"?>]
@PYaZ[<cfsSimulation] @PYaQ[xmlns=]@PYaB["http://www.cfs++.org"]@PYaZ[>]

 @PYaZ[<fileFormats]@PYaZ[>]
    @PYaZ[<output]@PYaZ[>]
      @PYaZ[<gid]@PYaZ[/>]
    @PYaZ[</output>]
    @PYaZ[<materialData] @PYaQ[file=]@PYaB["mat.xml"] @PYaQ[format=]@PYaB["xml"]@PYaZ[/>]
 @PYaZ[</fileFormats>]

 @PYaZ[<domain] @PYaQ[geometryType=]@PYaB["plane"]@PYaZ[>]
    @PYaZ[<regionList]@PYaZ[>]
      @PYaZ[<region] @PYaQ[name=]@PYaB["myCoil"] @PYaQ[material=]@PYaB["air"]@PYaZ[/>]
      @PYaZ[<region] @PYaQ[name=]@PYaB["iron"]   @PYaQ[material=]@PYaB["iron"]@PYaZ[/>]
      @PYaZ[<region] @PYaQ[name=]@PYaB["air"]    @PYaQ[material=]@PYaB["air"]@PYaZ[/>]
    @PYaZ[</regionList>]
    @PYaZ[<nodeList]@PYaZ[>]
      @PYaZ[<nodes] @PYaQ[name=]@PYaB["bound"]@PYaZ[/>]
    @PYaZ[</nodeList>]
 @PYaZ[</domain>]

 @PYaZ[<sequenceStep]@PYaZ[>]
    @PYaZ[<analysis]@PYaZ[>]
      @PYaZ[<static]@PYaZ[/>] 
    @PYaZ[</analysis>]
   @PYaZ[<pdeList]@PYaZ[>]
      @PYaZ[<magnetic]@PYaZ[>]
        @PYaZ[<regionList]@PYaZ[>]
          @PYaZ[<region] @PYaQ[name=]@PYaB["myCoil"] @PYaZ[/>]
          @PYaZ[<region] @PYaQ[name=]@PYaB["iron"]   @PYaZ[/>]
          @PYaZ[<region] @PYaQ[name=]@PYaB["air"]    @PYaZ[/>]
        @PYaZ[</regionList>]
        @PYaZ[<bcsAndLoads]@PYaZ[>]
          @PYaZ[<dirichletHom]   @PYaQ[name=]@PYaB["bound"]@PYaZ[/>]
        @PYaZ[</bcsAndLoads>]

        @PYaZ[<coils]@PYaZ[>]
          @PYaZ[<currentCoil2d] @PYaQ[name=]@PYaB["myCoil"]@PYaZ[>]
            @PYaE[<!--]@PYaE[ whole coil area: 100*5 mm^2 with 20 windings ]@PYaE[-->]
            @PYaZ[<windingCrossSection]@PYaZ[>]50e-6 @PYaZ[</windingCrossSection>]
            @PYaZ[<value]@PYaZ[>]       1            @PYaZ[</value>]
            @PYaZ[<id]@PYaZ[>]          1            @PYaZ[</id>]
          @PYaZ[</currentCoil2d>]
        @PYaZ[</coils>]

        @PYaZ[<storeResults]@PYaZ[>]
          @PYaZ[<elemResult] @PYaQ[type=]@PYaB["magFluxDensity"]@PYaZ[>]
            @PYaZ[<regionList]@PYaZ[>]
              @PYaZ[<region] @PYaQ[name=]@PYaB["myCoil"] @PYaZ[/>]
            @PYaZ[</regionList>]
          @PYaZ[</elemResult>]
          @PYaZ[<elemResult] @PYaQ[type=]@PYaB["magFluxDensity"]@PYaZ[>]
            @PYaZ[<regionList]@PYaZ[>]
              @PYaZ[<region] @PYaQ[name=]@PYaB["iron"]@PYaZ[/>]
            @PYaZ[</regionList>]
          @PYaZ[</elemResult>]
          @PYaZ[<elemResult] @PYaQ[type=]@PYaB["magFluxDensity"]@PYaZ[>]
            @PYaZ[<regionList]@PYaZ[>]
              @PYaZ[<region] @PYaQ[name=]@PYaB["air"]@PYaZ[/>]
            @PYaZ[</regionList>]
          @PYaZ[</elemResult>]
         @PYaZ[</storeResults>]
      @PYaZ[</magnetic>]
   @PYaZ[</pdeList>]
 @PYaZ[<sequenceStep]@PYaZ[>]
@PYaZ[</cfsSimulation>]
\end{Verbatim}

}

\newpage
\section{Heat Conduction}
\subsection{Results}
Until now, only the temperature, which is a nodal result, can be chosen for
output either for the whole domain or at named history nodes.

\begin{Verbatim}[commandchars=@\[\]]
@PYaZ[<storeResults]@PYaZ[>]
   @PYaZ[<nodeResult] @PYaQ[type=]@PYaB["heatTemperature"]@PYaZ[/>]
@PYaZ[</storeResults>]
\end{Verbatim}



\subsection{Example}

{\footnotesize
\begin{Verbatim}[commandchars=@\[\]]
@PYaN[<?xml version="1.0"?>]
@PYaZ[<cfsSimulation] @PYaQ[xmlns=]@PYaB["http://www.cfs++.org"]@PYaZ[>]

 @PYaZ[<fileFormats]@PYaZ[>]
    @PYaZ[<output]@PYaZ[>]
      @PYaZ[<gmv] @PYaQ[fixedGrid=]@PYaB["yes"] @PYaQ[binaryFormat=]@PYaB["no"]@PYaZ[/>]
    @PYaZ[</output>]
    @PYaZ[<materialData] @PYaQ[file=]@PYaB["mat.xml"] @PYaQ[format=]@PYaB["xml"]@PYaZ[/>]
 @PYaZ[</fileFormats>]

 @PYaZ[<domain] @PYaQ[geometryType=]@PYaB["plane"]@PYaZ[>]
    @PYaZ[<regionList]@PYaZ[>]
      @PYaZ[<region] @PYaQ[name=]@PYaB["slab"] @PYaQ[material=]@PYaB["steel"] @PYaZ[/>]
    @PYaZ[</regionList>]
 @PYaZ[</domain>]

 @PYaZ[<sequenceStep]@PYaZ[>]
   @PYaZ[<analysis]@PYaZ[>]
     @PYaZ[<transient]@PYaZ[>]
       @PYaZ[<numSteps]@PYaZ[>]     400 @PYaZ[</numSteps>]
       @PYaZ[<firstDt]@PYaZ[>]        2 @PYaZ[</firstDt>]
       @PYaZ[<stepSaveBeg]@PYaZ[>]    1 @PYaZ[</stepSaveBeg>]
       @PYaZ[<stepSaveEnd]@PYaZ[>]  400 @PYaZ[</stepSaveEnd>]
       @PYaZ[<stepSaveInc]@PYaZ[>]    1 @PYaZ[</stepSaveInc>]
       @PYaZ[<timeDataFile] @PYaQ[name=]@PYaB["sin80no.dat"]@PYaZ[/>]
     @PYaZ[</transient>] 
   @PYaZ[</analysis>]

   @PYaZ[<pdeList]@PYaZ[>]
     @PYaZ[<heatConduction]@PYaZ[>]
       @PYaZ[<regionList]@PYaZ[>]
         @PYaZ[<region] @PYaQ[name=]@PYaB["slab"]@PYaZ[/>]
       @PYaZ[</regionList>]
       @PYaZ[<bcsAndLoads]@PYaZ[>]
         @PYaZ[<dirichletInhom] @PYaQ[name=]@PYaB["bcleft"] @PYaQ[value=]@PYaB["0.0"]@PYaZ[/>]
         @PYaZ[<dirichletInhom] @PYaQ[name=]@PYaB["bcright"] @PYaQ[value=]@PYaB["100.0"] @PYaQ[dynamics=]@PYaB["sin80no.dat"]@PYaZ[/>]
       @PYaZ[</bcsAndLoads>]
       @PYaZ[<storeResults]@PYaZ[>]
         @PYaZ[<nodeResult] @PYaQ[type=]@PYaB["heatTemperature"]@PYaZ[>]
           @PYaZ[<nodeList]@PYaZ[>]
             @PYaZ[<nodes] @PYaQ[name=]@PYaB["historyNode"]@PYaZ[/>]
           @PYaZ[<nodeList]@PYaZ[>]
         @PYaZ[</nodeResult>]
       @PYaZ[</storeResults>]
     @PYaZ[</heatConduction>]
   @PYaZ[</pdeList>] 
 @PYaZ[<sequenceStep]@PYaZ[>]
@PYaZ[</cfsSimulation>]
\end{Verbatim}

}
